% ----------------------------------------
%
% LaTeX Article Template for CUED Reports
% Jon Sowman - 2009
% jon@hexoc.com
% 
% ----------------------------------------
% Set up the document class for an article
\documentclass[11pt]{article}

% Import the required packages for latex
\usepackage{appendix}

% This packages permits using $ \therefore $
\usepackage{amssymb}
\usepackage{graphicx}

% This package allows the use of $ \text{} $
\usepackage{amsmath}
\usepackage{savetrees}

% The document title and author
\title{SB3 - Datalogger\\Cambridge University Engineering Department} 
\author{David Turner \& Jon Sowman}

% Begin the document
\begin{document}
    \maketitle
	
% Insert the abstract for the document here
\begin{abstract}
    A compact, high speed logic analyser is designed with onboard buffering for
    the simultaneous analysis and logging of up to eight digital channels, plus
    an optional clock line for analysis of synchronous transfers. The supporting
    desktop application is also developed to configure and control the analyser,
    as well as to retrieve and post-process the recorded data and display it in
    a convenient format.
\end{abstract}

\section{Hardware Overview}
    A PIC18F4550 microcontroller on a breakout board provides the core of the
    datalogging hardware, clocked at 48MHz. The USB peripheral external hardware
    already exists on the board, along with hardware to bring the device out of
    reset into either the bootloader or user code.

    Additional hardware will be put in place to provide the following
    functionality:
    \begin{itemize}
    \item Monitor and log up to eight digital channels plus a clock.
    \item Store up to 1Mbit of data (that is to say, 128k samples if using all
    eight channels).
    \item Retrieve the samples from memory and transmit them to the desktop
    computer via USB interface for post processing, analysis and charting.
    \end{itemize}

\subsection{Filtering}
    Frontend antialiasing filters will be put in place to avoid frequency domain
    aliasing caused by signals at the input above the Nyquist frequency, which
    will be defined by the sampling frequency of the device. Simple first order
    low pass RC filters will suffice for this, with a $-3dB$ frequency placed
    just above the Nyquist frequency.

\subsection{Buffering}
    The datalogger will dump data to an SRAM buffer during the logging process,
    since the PIC does not have enough internal RAM to store a suitable number
    of samples to satisfy the specification. SRAM, whilst relatively expensive,
    also demonstrates extremely fast write speeds which is essential to keeping
    the sampling rate high.

    The hardware of choice is the AS6C1008 1Mbit SRAM IC from Alliance Memory, 
    which is divided into eight blocks. A 17-bit wide parallel interface is 
    used for byte-addressing, and a byte-wide parallel interface is provided 
    for data.

    The address lines will be directly attached to the PIC using the two
    byte-wide ports PORTB and PORTD, plus one additional line from another port.
    After the hardware anti-aliasing filters, the eight input channels are
    connected directly to the SRAM data interface via a 8-gate tri-state 
    buffer, controlled by the PIC. This buffer is enabled during the data
    acquisition phase, and subsequently disabled to detach the input channels
    from the memory.

\subsection{Data Retrieval}
    Retrieving data from the SRAM is achieved via the use of a
    parallel-in/serial-out shift register such as the 74HC165 series. The
    byte-wide parallel input to this device is attached to the SRAM data
    interface, and data is clocked out to the PIC before being packetised and
    transmitted to the desktop computer via the USB interface.

    A diagram of how to set up the apparatus is shown in figure \ref{fig:app}.
    
\subsection{Software Interface}
    The software interface will be constructed using LabWindows and will allow
    the initiation of datalogging and the viewing of captured data in several
    formats.  Asyncronous sample rate or synchronous mode, channels to
    capture, and number of samples to capture.  Once the full dataset has
    streamed from the PIC, it will be available for viewing in both timing diagram
    and listing formats.  There will be options to display listings in decimal, hex,
    binary, or ASCII.
    

    \begin{figure}
    \centering
    \includegraphics[height=5cm]{blank}
    \caption{Apparatus setup}
    \label{fig:app}
    \end{figure}
	
% Now print the appendix section
\appendix
\appendixpage
\addappheadtotoc
\section{Appendix 1}
			
\end{document}
